\chapter{Three questions about Spec, visually}\label{ch:refC-spec-zx}

\noindent\textit{The three questions that make $\Spec$ click --- drawn, then fully
explained. Each builds directly on Lesson 0001 (points $=$ primes, $V(I)$, $D(f)$,
generic/closed points, residue fields).}

\section{Draw me $\Spec\ZZ[x]$}\label{sec:zx-draw}

This is Mumford's \term{arithmetic surface} --- the friendliest 2-dimensional
spectrum. Because $\ZZ[x]$ is a UFD, its height-1 primes are exactly the principal
primes $(g)$ with $g$ irreducible, and there are only \textbf{two kinds} of
irreducible element --- so only two kinds of ``curve.''

\begin{figure}[h]
\centering
\resizebox{\textwidth}{!}{%
\begin{tikzpicture}[x=1pt,y=-1pt,scale=0.82,every node/.style={font=\small}]
  % frame
  \fill[paper2] (90,40) rectangle (650,340);
  \draw[rulec] (90,40) rectangle (650,340);
  % generic-point label
  \node[accent,font=\itshape\footnotesize,anchor=south] at (370,32)
    {$(0)$: generic point --- $\dim 2$, dense in all of $\Spec\ZZ[x]$};
  % vertical fibers V(p)
  \foreach \x in {170,270,370,470} \draw[accentB,line width=1.6pt] (\x,48)--(\x,332);
  \draw[gray,line width=1.2pt,dashed] (590,48)--(590,332);
  \node[inksoft] at (530,200) {$\cdots$};
  \node[inksoft,font=\itshape\scriptsize] at (590,62) {$\AA^1_{\QQ}$};
  % horizontal curves V(f)
  \draw[algc,line width=1.6pt] (96,110)--(644,110);
  \draw[algc,line width=1.6pt]
    (96,252) .. controls (150,227) and (210,277) .. (270,254)
             .. controls (330,232) and (410,277) .. (470,252)
             .. controls (530,232) and (600,274) .. (644,252);
  \node[algc,font=\footnotesize,anchor=south] at (126,108) {$V(x)$};
  \node[algc,font=\footnotesize,anchor=south] at (150,250) {$V(x^2+1)$};
  \node[inksoft,font=\itshape\scriptsize] at (600,300) {a 2-to-1 cover $\downarrow$};
  % closed points on V(x)
  \foreach \x in {170,270,370,470} \fill[accent] (\x,110) circle (3.6pt);
  \node[accent,font=\scriptsize,anchor=north] at (170,118) {$(2,x)$};
  \node[accent,font=\scriptsize,anchor=north] at (470,118) {$(7,x)$};
  % projection arrows
  \foreach \x in {170,270,370,470,590}
    \draw[gray,dashed,-{Stealth[length=4pt]}] (\x,340)--(\x,368);
  \node[inksoft,font=\itshape\footnotesize,anchor=west] at (606,358) {$\pi$};
  % base Spec Z
  \draw[inksoft,line width=1pt,-{Stealth[length=5pt]}] (90,375)--(660,375);
  \node[inksoft,font=\scriptsize,anchor=south west] at (92,373) {$\Spec\ZZ$};
  \foreach \x in {170,270,370,470} \fill[inksoft] (\x,375) circle (3pt);
  \draw[inksoft,line width=1pt] (590,375) circle (3pt);
  \node[inksoft,font=\scriptsize,anchor=north] at (170,383) {$(2)$};
  \node[inksoft,font=\scriptsize,anchor=north] at (270,383) {$(3)$};
  \node[inksoft,font=\scriptsize,anchor=north] at (370,383) {$(5)$};
  \node[inksoft,font=\scriptsize,anchor=north] at (470,383) {$(7)$};
  \node[inksoft,font=\scriptsize,anchor=north] at (590,383) {$(0)$};
\end{tikzpicture}}
\par\smallskip
{\footnotesize\color{inksoft}
\textcolor{accentB}{\rule[0.25ex]{1.6em}{1.5pt}}\ vertical $V(p)$ (a prime $p$, $\dim 1$)\quad
\textcolor{algc}{\rule[0.25ex]{1.6em}{1.5pt}}\ horizontal $V(f)$ (irreducible $f$, $\dim 1$)\quad
\textcolor{accent}{$\bullet$}\ closed point $(p,f)$ (maximal, $\dim 0$)}
\caption{$\Spec\ZZ[x]$ fibers over $\Spec\ZZ$ via $\pi$. The fiber over $(p)$ is the
vertical line $V(p)=\AA^1_{\FF_p}$; the fiber over the generic point $(0)$ is
$\AA^1_{\QQ}$.}
\end{figure}

The points sort cleanly by dimension:

\begin{center}
\begin{tabular}{@{}L{3.5cm}L{4.6cm}cl@{}}
\toprule
\textcolor{algc}{Prime of $\ZZ[x]$} & \textcolor{accentB}{Geometry} & \textbf{dim} & \textbf{residue field}\\
\midrule
\alg{$(0)$} & \geo{the one generic point, dense everywhere} & $2$ & $\QQ(x)$\\
\alg{$(p)$} & \geo{vertical line $V(p)=\AA^1_{\FF_p}$} & $1$ & $\FF_p(x)$\\
\alg{$(f)$, $f$ irred.\ primitive} & \geo{horizontal curve} & $1$ & a number field\\
\alg{$(p,f)$, $\bar f$ irred.\ mod $p$} & \geo{closed point (curve $\cap$ line)} & $0$ & finite $\FF_{p^{\deg\bar f}}$\\
\bottomrule
\end{tabular}
\end{center}

\begin{intubox}[The arithmetic hiding in the picture]
How a horizontal curve $V(f)$ meets a vertical line $V(p)$ is governed by how $f$
factors mod $p$. The curve $V(x^2+1)$ is literally $\Spec\ZZ[i]$, a 2-to-1 cover of
$\Spec\ZZ$ --- and its fiber over each prime is the splitting of that prime in the
Gaussian integers:
\end{intubox}

\begin{figure}[h]
\centering
\resizebox{\textwidth}{!}{%
\begin{tikzpicture}[x=1pt,y=-1pt,scale=0.82,every node/.style={font=\small}]
  \node[inksoft,font=\itshape\footnotesize] at (380,22)
    {fibers of $V(x^2+1)=\Spec\ZZ[i]$ over $\Spec\ZZ$};
  % connectors
  \draw[gray!50,dashed] (160,195)--(160,100);
  \draw[gray!50,dashed] (280,195)--(280,100);
  \draw[gray!50,dashed] (400,195)--(400,80);
  \draw[gray!50,dashed] (400,195)--(400,120);
  \draw[gray!50,dashed] (520,195)--(520,100);
  \draw[gray!50,dashed] (630,195)--(630,100);
  % fiber points
  \draw[accent,line width=1.6pt] (160,100) circle (4.4pt);
  \fill[accent] (160,100) circle (1.8pt);
  \fill[accentB] (280,100) circle (4pt);
  \fill[good] (400,80) circle (4pt);
  \fill[good] (400,120) circle (4pt);
  \fill[accentB] (520,100) circle (4pt);
  \fill[accentB] (630,100) circle (4pt);
  \node[accent,font=\scriptsize,anchor=south] at (160,76) {ramified};
  \node[accentB,font=\scriptsize,anchor=south] at (280,80) {inert $\cdot\ \FF_9$};
  \node[good,font=\scriptsize,anchor=south] at (400,60) {split $\cdot\ \FF_5\times 2$};
  \node[accentB,font=\scriptsize,anchor=south] at (520,80) {inert $\cdot\ \FF_{49}$};
  \node[accentB,font=\scriptsize,anchor=south] at (630,80) {inert $\cdot\ \FF_{121}$};
  \node[good,font=\itshape\scriptsize] at (430,138) {$5=(2+i)(2-i)$};
  % base
  \draw[inksoft,line width=1pt,-{Stealth[length=5pt]}] (80,195)--(690,195);
  \node[inksoft,font=\scriptsize,anchor=south west] at (82,193) {$\Spec\ZZ$};
  \foreach \x in {160,280,400,520,630} \fill[inksoft] (\x,195) circle (3pt);
  \node[inksoft,font=\scriptsize,anchor=north] at (160,203) {$(2)$};
  \node[inksoft,font=\scriptsize,anchor=north] at (280,203) {$(3)$};
  \node[inksoft,font=\scriptsize,anchor=north] at (400,203) {$(5)$};
  \node[inksoft,font=\scriptsize,anchor=north] at (520,203) {$(7)$};
  \node[inksoft,font=\scriptsize,anchor=north] at (630,203) {$(11)$};
\end{tikzpicture}}
\par\smallskip
{\footnotesize\color{inksoft}
\textcolor{good}{$\bullet$}\ split: $p\equiv1\ (4)$, two points\quad
\textcolor{accentB}{$\bullet$}\ inert: $p\equiv3\ (4)$, one point $\FF_{p^2}$\quad
\textcolor{accent}{$\circ$}\ ramified: $p=2$}
\caption{One irreducible curve, but a degree-2 cover: $p$ splits ($\equiv1$), is
inert ($\equiv3$), or ramifies ($=2$) --- the same data as factoring $x^2+1$ mod
$p$, i.e.\ primes in $\ZZ[i]$.}
\end{figure}

So the closed points \emph{are} the maximal ideals $(p,f)$, with
$\ZZ[x]/(p,f)=\FF_p[x]/(\bar f)$ a finite field --- and one wiggling curve quietly
encodes a chapter of number theory.

\section{Why isn't $\Spec(R)$ Hausdorff?}\label{sec:zx-hausdorff}

Hausdorff means any two distinct points sit in \emph{disjoint} open neighborhoods.
The Zariski topology has far too few opens for that --- for two reasons, both
visible already in $\Spec\ZZ$.

\begin{figure}[h]
\centering
\resizebox{\textwidth}{!}{%
\begin{tikzpicture}[x=1pt,y=-1pt,scale=0.78,every node/.style={font=\small}]
  \begin{scope}
    \node[good,font=\footnotesize] at (180,24) {$\RR$ --- Hausdorff $\checkmark$};
    \draw[gray,line width=1pt] (30,140)--(330,140);
    \draw[accentB,fill=accentB!10] (115,140) ellipse (48pt and 30pt);
    \draw[accentB,fill=accentB!10] (245,140) ellipse (48pt and 30pt);
    \fill[ink] (115,140) circle (3.2pt);
    \fill[ink] (245,140) circle (3.2pt);
    \node[ink,font=\footnotesize,anchor=north] at (115,148) {$a$};
    \node[ink,font=\footnotesize,anchor=north] at (245,148) {$b$};
    \node[inksoft,font=\itshape\scriptsize] at (180,200) {disjoint open intervals --- gap between};
    \node[good,font=\bfseries\footnotesize] at (180,224) {can separate any two points};
  \end{scope}
  \begin{scope}[xshift=380pt]
    \node[bad,font=\footnotesize] at (180,24) {$\Spec\ZZ$ --- not Hausdorff $\times$};
    % generic point smear
    \draw[accent!40,dashed,fill=accent!7] (190,140) ellipse (165pt and 16pt);
    \draw[gray,line width=1pt] (30,140)--(330,140);
    % overlapping opens
    \draw[bad,fill=bad!10] (135,140) ellipse (80pt and 34pt);
    \draw[bad,fill=bad!10] (225,140) ellipse (80pt and 34pt);
    \fill[accent] (120,140) circle (3.2pt);
    \fill[accent] (225,140) circle (3.2pt);
    \fill[inksoft] (300,140) circle (2.5pt);
    \node[accent,font=\footnotesize,anchor=south] at (120,128) {$(2)$};
    \node[accent,font=\footnotesize,anchor=south] at (225,128) {$(3)$};
    \node[accent,font=\itshape\scriptsize] at (46,166) {$(0)$};
    \node[inksoft,font=\itshape\scriptsize] at (180,200) {opens are cofinite --- they always overlap};
    \node[accent,font=\itshape\scriptsize] at (180,216) {\dots and the generic point $(0)$ is inside every open};
    \node[bad,font=\bfseries\footnotesize] at (180,236) {any two nonempty opens overlap};
  \end{scope}
\end{tikzpicture}}
\caption{On $\RR$ any two points are separated by disjoint opens; on $\Spec\ZZ$ every
nonempty open is cofinite, so any two of them overlap --- and the generic point lives
in all of them.}
\end{figure}

Pinning down the two reasons:
\begin{itemize}[leftmargin=1.4em,itemsep=3pt]
  \item \textbf{Points need not be closed.} Hausdorff $\Rightarrow$ $T_1$ (every
  point closed). But the generic point $(0)$ of a domain has closure the
  \emph{whole space} --- and $(0)\in D(f)$ for every $f\neq0$, so it sits inside
  \emph{every} nonempty open. You can never wall it off.
  \item \textbf{Nonempty opens always meet.} For an irreducible $\Spec(R)$ (e.g.\ $R$
  a domain), $D(f)\cap D(g)=D(fg)\neq\varnothing$ whenever $f,g\neq0$. In $\Spec\ZZ$
  the only proper closed sets are \emph{finite} sets of primes, so every nonempty
  open is cofinite, and two cofinite sets always intersect.
\end{itemize}

\begin{primarybox}
\centering
$\Spec(R)$ is Hausdorff $\iff$ it is $T_1$ $\iff$ every prime is maximal
$\iff$ $\dim R = 0$.
\end{primarybox}

So the only Hausdorff spectra are the dimension-0 ones: a field gives a single
point; an Artinian ring gives a finite discrete set; a Boolean ring gives a
Cantor-like Stone space. Everything with actual geometric dimension is \emph{not}
Hausdorff --- and that's by design. The Zariski topology is deliberately coarse;
the geometric richness is moved into the structure sheaf $\O_X$ instead. Generic
points aren't a bug --- they let you speak of ``the behavior at the generic
point,'' i.e.\ the function field.

\section{What's the residue field at $(p)$?}\label{sec:zx-residue}

An element $f\in R$ behaves like a \emph{function} on $\Spec(R)$. Its value at a
point $\fp$ lives in the \textbf{residue field}

\begin{primarybox}
\centering
$\displaystyle \res(\fp)\;=\;\Frac(R/\fp)\;=\;R_{\fp}/\fp R_{\fp},
\qquad\text{and}\qquad f(\fp):=\big[f\big]\in\res(\fp).$
\end{primarybox}

\begin{figure}[h]
\centering
\resizebox{\textwidth}{!}{%
\begin{tikzpicture}[x=1pt,y=-1pt,scale=0.82,every node/.style={font=\small}]
  \node[inksoft,font=\footnotesize] at (380,24)
    {``evaluate the function $n$ at the point $(p)$'' $=$ ``reduce $n$ mod $p$''};
  % source box
  \draw[algc,line width=1pt,fill=paper2,rounded corners=6pt] (60,70) rectangle (240,150);
  \node[algc,font=\large] at (150,96) {$n\in\ZZ$};
  \node[inksoft,font=\itshape\scriptsize] at (150,124) {a function on $\Spec\ZZ$};
  % arrow
  \draw[accentB,line width=1.1pt,-{Stealth[length=5pt]}] (250,110)--(378,110);
  \node[accentB,font=\scriptsize,anchor=south] at (315,104) {at $(p)$: $n\mapsto n\bmod p$};
  % target circle
  \draw[accentB,line width=1pt,fill=accentB!8] (470,110) circle (52pt);
  \node[accentB,font=\large] at (470,104) {$\FF_p$};
  \node[inksoft,font=\itshape\scriptsize] at (470,128) {$=\res((p))$};
  % example box
  \draw[good,line width=1pt,fill=paper2,rounded corners=6pt] (556,74) rectangle (726,146);
  \node[good,font=\footnotesize] at (641,100) {$100$ at $(7)$};
  \node[ink,font=\footnotesize] at (641,124) {$=2\quad(100=14\cdot7+2)$};
  % footnote line
  \node[inksoft,font=\itshape\scriptsize] at (470,196)
    {since $\ZZ/(p)$ is already a field, $\res((p))=\FF_p$; at the generic point $\res((0))=\QQ$};
\end{tikzpicture}}
\caption{The residue field at $(p)\in\Spec\ZZ$ is $\FF_p$ --- so ``reducing $n$ mod
$p$'' is \emph{literally} evaluating the function $n$ at the point $(p)$.}
\end{figure}

\begin{intubox}[Try it]
Evaluate the integer $n$ at the point $(p)$. For example, $100 = 14\cdot 7 + 2$, so
the function $100$ evaluates to $2$ at the point $(7)$; the value lives in the residue
field $\res((7))=\FF_7$. The point $(N)$ is a genuine point with residue field
$\FF_N$ only when $N$ is prime --- for composite $N$, $\ZZ/(N)$ is not a field.
\end{intubox}

The pattern generalizes: \textbf{closed points carry small (finite) residue fields;
generic points carry big transcendental ones.} The local ring $R_{\fp}$ reveals the
dimension through its residue field. Reading it off the arithmetic surface from
\hyperref[sec:zx-draw]{Q1}:

\begin{center}
\begin{tabular}{@{}lll@{}}
\toprule
\textbf{Point} & $\res(\fp)$ & \textbf{flavor}\\
\midrule
$(0)$ --- generic, $\dim 2$ & $\QQ(x)$ & function field, transcendental\\
$(p)$ --- vertical line & $\FF_p(x)$ & function field over $\FF_p$\\
$(f)$ --- horizontal curve & $\Frac(\ZZ[x]/(f))$ & a number field\\
$(p,f)$ --- closed point & $\FF_{p^{\deg\bar f}}$ & a finite field\\
\bottomrule
\end{tabular}
\end{center}

And here Q1 and Q3 fuse: at the closed point $(p,x^2+1)$ over an \emph{inert} prime
the residue field is $\FF_{p^2}$, while over a \emph{split} prime the two points each
have residue field $\FF_p$ --- the residue field \emph{is} the splitting type.

\begin{intubox}[Want to go further?]
Ask me: ``draw $\Spec k[x,y]$,'' ``what does the local ring $R_{(p)}$ look like for
$\ZZ$?'', or ``show the structure sheaf's value on a $D(f)$ here.'' I'm your teacher
between lessons.
\end{intubox}
